\documentclass[12pt,a4paper,twoside]{report}
\usepackage[utf-8]{inputenc}
\usepackage[vietnamese]{babel}
\usepackage[margin=1in]{geometry}
\usepackage{amsmath}
\usepackage{amssymb}
\usepackage{graphicx}
\usepackage{hyperref}
\usepackage{float}
\usepackage{array}
\usepackage{booktabs}
\usepackage{multirow}
\usepackage{longtable}
\usepackage{xcolor}
\usepackage{textcomp}

\hypersetup{
    colorlinks=true,
    linkcolor=blue,
    urlcolor=blue,
    citecolor=blue
}

\renewcommand{\arraystretch}{1.5}

\title{Network Intrusion Detection using Reinforcement Learning\\
An RL-based Defense Agent with Adaptive Policy Learning}
\author{Group Member Names}
\date{Hanoi, 2026}

\begin{document}

\maketitle

\newpage

\chapter*{ABSTRACT}

\section*{Network Intrusion Detection using Reinforcement Learning: An RL-based Defense Agent with Adaptive Policy Learning}

\noindent
\textbf{Background:} Traditional network defense relies on static rules and signature-based detection, struggling to adapt to automated, evolving cyber attacks. This work addresses the need for autonomous defense systems using Reinforcement Learning (RL) to learn optimal defensive policies.

\vspace{0.5cm}
\noindent
\textbf{Methods:} We developed an RL-based Network Intrusion Detection System agent using Proximal Policy Optimization trained in simulation with 5 attack types (SYN Flood, Port Scan, Brute-Force, SQL Injection, XSS). The agent observes a 20-dimensional feature vector combining network-level metrics (packet rate, inter-arrival time, port distribution) and HTTP payload features (SQL/XSS detection via OWASP CRS). It selects from three defensive actions: Block, RateLimit, or Redirect-to-Honeypot. Training proceeded for ~350,000 timesteps with evaluation every 12,000 steps.

\vspace{0.5cm}
\noindent
\textbf{Results:} The RL agent achieved \textbf{91.6\% detection rate}, outperforming Static Rules (74.3\%) and Random Forest (87.9\%). Per-attack: SYN Flood 99.5\%, Port Scan 67.4\%, SQLi 73.2\%, XSS 81.4\%, Brute-Force 88.5\%. The policy balanced security and availability: 93.2\% legitimate traffic allowed (FPR 6.8\%), 96.4\% SYN Flood mitigation. Training converged smoothly with +475\% reward improvement, safe policy updates (KL-divergence: 0.012$\to$0.003). Ablation study showed the HTTP payload feature group as the dominant contributor ($-$18.4~pp), with payload normalization ($-$12.8~pp) exceeding CRS rules ($-$10.2~pp).

\vspace{0.5cm}
\noindent
\textbf{Conclusions:} This work demonstrates RL's feasibility for autonomous network defense. The agent learns stable, differentiated policies across diverse attacks while balancing mitigation with availability. A hybrid architecture combining static rate-limiting with RL for application-layer analysis is necessary to address system latency constraints. Future work includes multi-agent coordination and simulation-to-reality validation.

\vspace{0.7cm}
\noindent
\textbf{Keywords:} Reinforcement Learning, Network Intrusion Detection, Proximal Policy Optimization, Cybersecurity Defense, Feature Engineering, Payload Normalization, OWASP ModSecurity, Autonomous Defense, Policy Learning, Network Analysis

\newpage

\tableofcontents

\newpage

\chapter*{LIST OF FIGURES}

\setlength{\parindent}{0pt}
\vspace{0.5cm}

\noindent
Figure 4.1: Đường cong học tập \texttt{eval/mean\_reward} và \texttt{train/entropy\_loss} theo số bước huấn luyện \dotfill~Page~\pageref{fig:learning_curves}

\newpage

\chapter*{LIST OF TABLES}

\setlength{\parindent}{0pt}
\vspace{0.5cm}

\noindent
Table 3.1: Danh sách 20 đặc trưng trong vector quan sát của tác tử RL \\
Table 3.2: Ma trận Phủ nhận --- Đặc trưng $\times$ Loại Tấn công \\
Table 4.1: Kết quả phân loại của từng nhóm đặc trưng \\
Table 4.2: Chỉ số chẩn đoán PPO trong quá trình huấn luyện \\
Table 4.2b: Hiệu suất huấn luyện --- Default SB3 vs. Đã tinh chỉnh \\
Table 4.3: Kết quả phân loại hành động của Tác tử Phòng thủ RL \\
Table 4.4: Nghiên cứu Ablation --- Ảnh hưởng khi loại bỏ từng thành phần \\
Table 4.5: So sánh hiệu suất tác tử RL với các phương pháp cơ sở

\newpage

\chapter*{ABBREVIATIONS}

\begin{longtable}[c]{p{3cm}p{12cm}}
\toprule
\textbf{Viết tắt} & \textbf{Ý nghĩa đầy đủ} \\
\midrule
AI & Artificial Intelligence (Trí tuệ Nhân tạo) \\
API & Application Programming Interface \\
CAPTCHA & Completely Automated Public Turing test to tell Computers and Humans Apart \\
CRS & Core Rule Set (Tập Quy tắc Cốt lõi) \\
CSV & Comma-Separated Values \\
DDoS & Distributed Denial of Service (Tấn công Từ chối Dịch vụ Phân tán) \\
DoS & Denial of Service (Tấn công Từ chối Dịch vụ) \\
RL & Reinforcement Learning (Học tăng cường) \\
FPR & False Positive Rate (Tỷ lệ Dương tính Giả) \\
HTTP & Hypertext Transfer Protocol \\
HTTPS & Hypertext Transfer Protocol Secure \\
IDS & Intrusion Detection System (Hệ thống Phát hiện Xâm nhập) \\
NIDS & Network Intrusion Detection System (Hệ thống Phát hiện Xâm nhập Mạng) \\
KL & Kullback-Leibler (Divergence) \\
MDP & Markov Decision Process (Quá trình Quyết định Markov) \\
MIME & Multipurpose Internet Mail Extensions \\
ML & Machine Learning (Học Máy) \\
MTU & Maximum Transmission Unit \\
NFKC & Compatibility Decomposition, followed by Canonical Composition \\
NLP & Natural Language Processing \\
OWASP & Open Web Application Security Project \\
PCAP & Packet Capture Format \\
PL & Paranoia Level \\
PPO & Proximal Policy Optimization (Tối ưu hóa Chính sách Xấp xỉ) \\
RQ & Research Question (Câu hỏi Nghiên cứu) \\
RFC & Request for Comments \\
SB3 & Stable-Baselines3 \\
SDN & Software-Defined Networking (Mạng Điều khiển bằng Phần mềm) \\
SIEM & Security Information and Event Management \\
SLA & Service Level Agreement \\
SOC & Security Operations Center (Trung tâm Điều hành An ninh) \\
SQL & Structured Query Language \\
SQLi & SQL Injection (Chèn mã SQL) \\
SSH & Secure Shell \\
SSL/TLS & Secure Sockets Layer / Transport Layer Security \\
SYN & Synchronization flag (TCP) \\
TCP & Transmission Control Protocol \\
TLS & Transport Layer Security \\
TTPs & Tactics, Techniques, and Procedures \\
UDP & User Datagram Protocol \\
VM & Virtual Machine (Máy ảo) \\
XAI & Explainable Artificial Intelligence \\
XML & Extensible Markup Language \\
XSS & Cross-Site Scripting (Tập lệnh Trên nhiều Trang) \\
Zero-day & Previously unknown software vulnerability \\
\bottomrule
\end{longtable}

\newpage

% ========== CHAPTER 1: INTRODUCTION ==========
\chapter{GIỚI THIỆU}

\section{Bối cảnh}

Bối cảnh an ninh mạng đã trải qua một sự chuyển đổi đáng kể trong thập kỷ qua, được thúc đẩy bởi sự tăng trưởng nhanh chóng của cơ sở hạ tầng kỹ thuật số, điện toán đám mây và các mạng lưới kết nối quy mô lớn. Song song với những phát triển này, các mối đe dọa mạng ngày càng trở nên tự động hóa, thích ứng và tinh vi hơn, đặt ra những thách thức nghiêm trọng đối với các cơ chế bảo mật truyền thống.

\subsection{Sự leo thang của các mối đe dọa mạng và Tự động hóa tấn công}

Các cuộc tấn công mạng hiện đại không còn chủ yếu là thủ công hay mang tính cơ hội. Thay vào đó, kẻ thù ngày càng dựa vào các công cụ tự động để thực hiện trinh sát quy mô lớn, quét lỗ hổng, nhồi nhét thông tin xác thực (credential stuffing) và khai thác ở tốc độ máy.

\subsection{Tác động kinh tế của vi phạm dữ liệu}

\subsection{Sự phát triển của cơ sở hạ tầng mạng và Thách thức vận hành}

\subsection{Nhu cầu về Hệ thống phòng thủ tự chủ và thích ứng}

\section{Phát biểu vấn đề}

\section{Mục tiêu nghiên cứu}

\section{Ý nghĩa của nghiên cứu}

\section{Phạm vi và Hạn chế}

\subsection{Phạm vi nghiên cứu}

\subsection{Hạn chế của Phương pháp}

% ========== CHAPTER 2: LITERATURE REVIEW ==========
\chapter{TỔNG QUAN TÀI LIỆU}

\textit{[Content from chapter\_full\_vietnamese.md Chapter 2]}

\section{Tổng quan Hệ thống Phát hiện Xâm nhập (IDS)}

\section{Quản lý Luồng Mạng (Network Flow Management)}

\section{Kỹ thuật Trích xuất Đặc trưng và Quy chuẩn}

\section{Học tăng cường (Reinforcement Learning) cho Bảo mật Mạng}

\section{Học Chính sách và Thuật toán Tối ưu hóa}

\section{Quy tắc ModSecurity OWASP CRS}

\section{So sánh với Công trình Liên quan}

% ========== CHAPTER 3: METHODOLOGY ==========
\chapter{PHƯƠNG PHÁP NGHIÊN CỨU}

\textit{[Content from chapter\_full\_vietnamese.md Chapter 3]}

\section{Thiết kế Nghiên cứu}

\subsection{Mô hình Mối đe dọa}

\subsection{Khung Phương pháp Nghiên cứu}

\subsection{Kiến trúc Tổng thể Hệ thống}

\subsection{Hình thức hóa Bài toán MDP}

\subsection{Topology Mạng Thực nghiệm}

\subsection{Thiết kế Tác tử Phòng thủ RL}

\subsection{Pipeline Suy diễn và Thực thi Thời gian Thực}

\section{Phương pháp Thu thập Dữ liệu}

\subsection{Bắt và Phân tách Gói tin}

\subsection{Quản lý Phiên kết nối}

\subsection{Cửa sổ Trượt 1 Giây}

\subsection{Nguồn Dữ liệu}

\section{Kỹ thuật Phân tích Dữ liệu}

\subsection{Tổng quan Kỹ nghệ Đặc trưng}

\subsection{Đặc trưng Mạng (F1--F11)}

\subsection{Đặc trưng SQL Injection (F12--F17)}

\subsection{Đặc trưng Cross-Site Scripting (F18--F20)}

\subsection{Pipeline Chuẩn hóa Payload}

\subsection{Ma trận Phủ nhận Tấn công}

\subsection{Chuẩn hóa Vector Đặc trưng}

\section{Hạn chế của Phương pháp}

% ========== CHAPTER 4: RESULTS ==========
\chapter{KẾT QUẢ THỰC NGHIỆM}

\textit{[Content from chapter\_full\_vietnamese.md Chapter 4]}

\section{Giới thiệu}

\section{Trình bày Dữ liệu}

\subsection{Kết quả Pipeline Chuẩn hóa Payload}

\subsection{Phân tích CRS Paranoia Level}

\subsection{Hiệu suất Phát hiện theo Đặc trưng}

\subsection{Trọng số Đặc trưng và Phân cụm}

\subsection{Đường cong Học tập của Tác tử RL}

\subsection{Chỉ số Chẩn đoán PPO}

\label{fig:learning_curves}
\subsection{So sánh Cấu hình Mặc định SB3 và Đã Tinh chỉnh}

\subsection{Kết quả Đánh giá Cuối cùng}

\section{Phân tích Kết quả}

\subsection{Đóng góp Thành phần (Nghiên cứu Ablation)}

\subsection{Hiệu năng Phản ứng Thời gian Thực}

\subsection{Phân tích Hành vi Chính sách}

\section{Diễn giải Kết quả}

\subsection{Đánh giá Tổng thể}

\subsection{Xác minh Kịch bản Thực tế}

\subsection{Phát hiện Kỹ thuật Chính}

\section{So sánh với Tài liệu}

\section{Hàm ý của Kết quả}

\subsection{Khi nào Nên Ưu tiên RL thay vì Luật Tĩnh}

\subsection{Yêu cầu Hạ tầng Tối thiểu}

\subsection{Đánh đổi Độ trễ so với Độ chính xác}

\subsection{Honeypot như Lợi thế Chiến lược}

\subsection{Khả năng Tái sử dụng Chính sách}

\subsection{Hướng Phát triển Tiếp theo}

% ========== CHAPTER 5: DISCUSSION ==========
\chapter{THẢO LUẬN}

\section{Phát biểu lại Câu hỏi Nghiên cứu và Mục tiêu}

Dự án này được dẫn dắt bởi ba câu hỏi nghiên cứu chính, mỗi câu hỏi đều nhằm xác định một khía cạnh khác nhau của tính khả thi và hiệu quả của phòng thủ mạng dựa trên RL.

\textbf{RQ1 --- Hiệu quả so sánh:} Tác tử học tăng cường có vượt trội hơn so với các cơ chế phòng thủ dựa trên quy tắc truyền thống (luật tĩnh, bộ phân loại máy học) trong việc phát hiện và giảm thiểu các cuộc tấn công mạng tự động và thích ứng trong môi trường mô phỏng?

\textbf{RQ2 --- Đánh đổi vận hành:} Làm thế nào mà việc thực hiện các hành động phòng thủ tự chủ (chặn, giới hạn tốc độ, chuyển hướng) ảnh hưởng đến sự cân bằng giữa việc giảm thiểu các cuộc tấn công thành công và tính sẵn sàng của dịch vụ hợp lệ?

\textbf{RQ3 --- Tính ổn định của chính sách:} Liệu tác tử dựa trên RL có thể xây dựng và điều chỉnh một chính sách phòng thủ ổn định, duy trì hiệu quả trên các vector tấn công đa dạng và không ngừng phát triển mà không cần cấu hình lại thủ công?

\section{Tóm tắt Phát hiện Chính và Diễn giải trong Bối cảnh Nghiên cứu}

\subsection{RQ1 --- Hiệu quả so sánh của Tác tử RL}

Kết quả thực nghiệm cho thấy tác tử RL đạt tỷ lệ phát hiện \textbf{91,6\%} trên tập hợp các cuộc tấn công đa dạng, vượt trội hơn so với hai baseline:
\begin{itemize}
\item \textbf{Luật Tĩnh (Static Rules):} 74,3\% --- độ lệch +17,3 pp
\item \textbf{Random Forest (RF) trên vector đặc trưng 20 chiều:} 87,9\% --- độ lệch +3,7 pp
\end{itemize}

\subsection{RQ2 --- Đánh đổi vận hành}

Phòng thủ mạng là bài toán đánh đổi vốn có: chúng ta muốn chặn tất cả các cuộc tấn công nhưng không muốn gây ảnh hưởng đến người dùng hợp lệ.

\textbf{Tính sẵn sàng của dịch vụ:}
\begin{itemize}
\item \textbf{Legitimate traffic allowed:} 93,2\%
\item \textbf{False Positive Rate (FPR):} 6,8\%
\end{itemize}

\subsection{RQ3 --- Tính ổn định của Chính sách}

Dữ liệu huấn luyện xác nhận rằng tác tử duy trì \textbf{ổn định policy cao} trên 5 loại tấn công đa dạng.

\subsection{Phân tích chi tiết: Các Nhân tố Chi phối}

\subsubsection{Nhân tố Chi phối: Nhóm Đặc trưng HTTP Payload (F12--F20)}

Ablation study chi tiết cho thấy:
\begin{itemize}
\item \textbf{Nhóm F12--F20:} $-$18,4 pp (largest single impact)
\item \textbf{PayloadNormalizer:} $-$12,8 pp (lớn hơn CRS)
\item \textbf{CRS rules (F13, F18):} $-$10,2 pp
\end{itemize}

PayloadNormalizer ($-$12,8 pp) $>$ CRS ($-$10,2 pp) từ khi xét single component.

\subsubsection{Tinh chỉnh Siêu tham số}

\begin{table}[h]
\centering
\begin{tabular}{|c|c|c|c|}
\hline
\textbf{Cấu hình} & \textbf{Reward Cuối} & \textbf{Hội tụ} & \textbf{Wall-clock} \\
\hline
Default SB3 & $\sim$2,36 & $\sim$80k & 786s \\
\textbf{Tuned} & \textbf{$\sim$2,28} & \textbf{$\sim$220k} & \textbf{197s} \\
\hline
\end{tabular}
\end{table}

\subsection{Hạn chế và Cảnh báo Quan trọng}

Năm hạn chế lớn cần ghi nhận:

\subsubsection{Khoảng cách Mô phỏng (Simulation Gap)}

Toàn bộ thử nghiệm diễn ra trong \textbf{Containernet}, một container ảo hóa của mạng linux-based.

\subsubsection{Phụ thuộc vào Hàm Phần thưởng (Reward Function Dependency)}

Phát hiện 91,6\% được đưa ra dưới \textbf{một cụ thể (α, β, γ)} --- nó không phải ``mức phát hiện tối ưu'' mà là ``mức phát hiện được học dưới design choice này.''

\subsubsection{Phân tích Per-window che giấu Một phần Điểm yếu}

Per-window analysis cho thấy 26,3\% recall cho SQLi, nhưng 75,4\% của windows không chứa payload.

\subsubsection{Ràng buộc HTTPS (HTTPS Constraint)}

20 đặc trưng dựa vào \textbf{plaintext inspection} của HTTP/TCP. Khi dữ liệu được mã hóa (TLS 1.3), F12--F20 không thể tính toán được.

\subsubsection{Khả năng Mở rộng Ngang (Horizontal Scalability)}

Dự án này training một agent duy nhất cho toàn bộ mạng.

% ========== CHAPTER 6: CONCLUSION ==========
\chapter{KẾT LUẬN VÀ HƯỚNG PHÁT TRIỂN}

\section{Kết Luận}

\subsection{Vấn đề Nghiên cứu Đã Được Xác định}

Chương 1 định rõ một nhu cầu: hệ thống phòng thủ mạng hiện đại là \textbf{phản ứng, tĩnh, và overload bởi alert fatigue}.

\subsection{Dự án Đã Chứng minh Tính Khả thi}

Kết quả thực nghiệm cho thấy \textbf{câu trả lời là có}, với các chứng cứ cụ thể:

\subsubsection{RL Hiệu quả So sánh (RQ1)}

\begin{itemize}
\item \textbf{RL Detection Rate: 91,6\%} --- so với Static Rules (74,3\%, $-$17,3 pp) và Random Forest (87,9\%, $-$3,7 pp)
\item HTTP payload feature group (F12--F20) với emphasis trên PayloadNormalizer đóng vai trò then chốt
\end{itemize}

\subsubsection{RL Đánh đổi Vận hành (RQ2)}

\begin{itemize}
\item \textbf{Legitimate traffic allowed: 93,2\%} (FPR 6,8\%)
\item \textbf{Attack mitigation: $>$96\% DoS, 100\% SQLi session-level}
\item Honeypot redirection strategy cung cấp dual-benefit: mitigation + threat intelligence
\end{itemize}

\subsubsection{RL Tính Ổn định Chính sách (RQ3)}

\begin{itemize}
\item Training curve smooth (+475\% reward), no oscillation
\item approx\_kl safely converges (0,012 $\to$ 0,003)
\item Agent learns differentiated responses per attack type
\end{itemize}

\subsection{Ba Đóng góp Chính của Dự án}

\begin{enumerate}
\item \textbf{Technical Contribution:} End-to-end RL NIDS integration. HTTP payload group (F12--F20) is dominant, not CRS alone.

\item \textbf{Methodological Contribution:} Rigorous evaluation framework: per-window vs per-session analysis, ablation metrics, FPR trade-offs.

\item \textbf{Operational Contribution:} Deployment feasibility: 1.016 ms latency, 6.8\% FPR, multi-stage actions, SIEM-ready.
\end{enumerate}

\section{Hướng Phát triển Tiếp Theo}

\subsection{Mở rộng Không gian Hành động (Action Space Expansion)}

\textbf{Hiện tại:} Agent chọn giữa 3 hành động (Block, RateLimit, Redirect-to-Honeypot)

\textbf{Đề xuất:}
\begin{itemize}
\item Per-source-IP rate limiting
\item Per-request-type routing
\item Dynamic honeypot spawning
\item Action masking
\end{itemize}

\textbf{Effort:} Moderate (2--3 weeks)

\subsection{Phòng thủ Phân tán và Đa tác tử (Multi-Agent Coordination)}

\textbf{Hiện tại:} Single agent monitors entire network

\textbf{Đề xuất:}
\begin{itemize}
\item Agent per network segment
\item Distributed policy coordination (centralized or decentralized)
\item Communication protocol for threat alerts
\end{itemize}

\textbf{Effort:} Very High (4--6 months)

\subsection{Robustness đối với Adversarial Evasion}

\textbf{Hiện tại:} Attack evaluation is ``cooperative''

\textbf{Đề xuất:}
\begin{itemize}
\item Adversarial attacker model with feature knowledge
\item Evasion tactics: normalization bypass, timing variation, polyglot payloads
\item Adversarially train policy against adaptive attacker
\end{itemize}

\textbf{Effort:} Very High (5--6 months)

\subsection{Online Retraining và Concept Drift Handling}

\textbf{Hiện tại:} Policy fixed after training

\textbf{Đề xuất:}
\begin{itemize}
\item Drift detection monitoring
\item Incremental retraining with continual learning
\item A/B testing on shadow traffic
\end{itemize}

\textbf{Effort:} Moderate--High (2--3 weeks)

\subsection{Transparent Decision Logs và Explainability}

\textbf{Hiện tại:} Policy is black-box

\textbf{Đề xuất:}
\begin{itemize}
\item Attention visualization (SHAP, LIME per-inference)
\item Decision logs with feature explanation
\item Policy audit trail for regulatory compliance
\end{itemize}

\textbf{Effort:} Moderate (2--3 weeks)

\subsection{Simulation-to-Reality Gap Validation (Sim2Real Transfer)}

\textbf{Hiện tại:} Policy trained on Containernet, real-world unknown

\textbf{Phase 1: Distribution Shift Analysis}
\begin{itemize}
\item Collect real network PCAP
\item Compare F1--F20 statistics vs. simulated
\item Identify largest distribution shifts
\end{itemize}

\textbf{Phase 2: Policy Transfer + Evaluation}
\begin{itemize}
\item Deploy trained policy on real traffic (shadow mode)
\item Measure detection rate degradation
\end{itemize}

\textbf{Phase 3: Sim2Real Fine-tuning}
\begin{itemize}
\item Fine-tune policy on real data
\item Restore detection to near 91\%
\end{itemize}

\textbf{Effort:} Very High (2--3 months)

\subsection{Tóm tắt và Tầm nhìn Tương lai}

Dự án này đã chứng minh rằng RL là một \textbf{feasible, practical nền tảng} cho hệ thống phòng thủ mạng tự chủ. Ba câu hỏi nghiên cứu được trả lời:
\begin{itemize}
\item \textbf{RQ1:} RL outperforms rule-based (91.6\% vs. 74.3\%)
\item \textbf{RQ2:} RL balances security $\land$ availability (93.2\% legit allowed, 96\%+ attack mitigation)
\item \textbf{RQ3:} RL maintains stable policy across diverse attacks
\end{itemize}

Sự kết hợp của \textbf{academic rigor} (ablation studies, distributed learning theory) và \textbf{operational pragmatism} (honeypot escalation, FPR management, SIEM integration) là chìa khóa để RL defense không chỉ là proof-of-concept mà là \textbf{deployed reality} trong các tổ chức sắp tới.

% ========== REFERENCES ==========
\begin{thebibliography}{99}

\bibitem{puterman1994} Puterman, M. L. (1994). \textit{Markov Decision Processes: Discrete Stochastic Dynamic Programming}. John Wiley \& Sons.

\bibitem{schulman2017} Schulman, J., Wolski, F., Dhariwal, P., Radford, A., \& Klimov, O. (2017). Proximal Policy Optimization Algorithms. \textit{arXiv preprint arXiv:1707.06347}.

\bibitem{raffin2021} Raffin, A., Hill, A., Gleave, A., Kanervisto, A., Ernestus, M., \& Dormann, N. (2021). Stable-Baselines3: Reliable Reinforcement Learning Implementations. \textit{Journal of Machine Learning Research}, \textbf{22}(268), 1--8.

\bibitem{peuster2016} Peuster, M., Karl, H., \& van Rossem, S. (2016). MeDICINE: Rapid Prototyping of Production-Ready Network Services in Multi-PoP Environments. In \textit{Proceedings of the IEEE Conference on Network Function Virtualization and Software Defined Networks (NFV-SDN 2016)} (pp. 148--153).

\bibitem{lashkari2024} Lashkari, A. H., et al. (2024). \textit{LSNM2024: A Labeled Network Traffic Dataset for Intrusion Detection}. University of New Brunswick.

\bibitem{sharafaldin2018} Sharafaldin, I., Lashkari, A. H., \& Ghorbani, A. A. (2018). Toward Generating a New Intrusion Detection Dataset and Intrusion Traffic Characterization. In \textit{Proceedings of the 4th International Conference on Information Systems Security and Privacy (ICISSP 2018)} (pp. 108--116).

\bibitem{postel1981} Postel, J. (1981). \textit{Transmission Control Protocol}. RFC 793. IETF.

\bibitem{owasp2023} OWASP Foundation. (2023). \textit{OWASP ModSecurity Core Rule Set (CRS) v4.0}. https://coreruleset.org/

\bibitem{spett2005} Spett, K. (2005). \textit{SQL Injection: Are Your Web Applications Vulnerable?} SPI Dynamics White Paper.

\bibitem{kruegel2003} Kruegel, C., \& Vigna, G. (2003). Anomaly Detection of Web-Based Attacks. In \textit{Proceedings of the 10th ACM CCS} (pp. 251--261).

\bibitem{towers2023} Towers, M., et al. (2023). Gymnasium: A Standard Interface for Reinforcement Learning Environments. \textit{arXiv preprint arXiv:2407.17032}.

\bibitem{sutton2018} Sutton, R. S., \& Barto, A. G. (2018). \textit{Reinforcement Learning: An Introduction} (2nd ed.). MIT Press.

\bibitem{mnih2015} Mnih, V., et al. (2015). Human-Level Control through Deep Reinforcement Learning. \textit{Nature}, \textbf{518}(7540), 529--533.

\bibitem{owasp2021} OWASP Foundation. (2021). \textit{OWASP Top Ten}. https://owasp.org/www-project-top-ten/

\bibitem{shiravi2012} Shiravi, A., Shiravi, H., Tavallaee, M., \& Ghorbani, A. A. (2012). Toward Developing a Systematic Approach to Generate Benchmark Datasets for Intrusion Detection. \textit{Computers \& Security}, \textbf{31}(3), 357--374.

\bibitem{moustafa2015} Moustafa, N., \& Slay, J. (2015). UNSW-NB15: A Comprehensive Dataset for Network Intrusion Detection Systems. In \textit{Proceedings of MilCIS 2015}.

\bibitem{ring2019} Ring, M., et al. (2019). A Survey of Network-Based Intrusion Detection Data Sets. \textit{Computers \& Security}, \textbf{86}, 147--167.

\bibitem{modsecurity2023} ModSecurity. (2023). \textit{ModSecurity Web Application Firewall Documentation}. https://modsecurity.org/

\bibitem{fielding2014} Fielding, R., \& Reschke, J. (2014). \textit{Hypertext Transfer Protocol (HTTP/1.1): Message Syntax and Routing}. RFC 7230. IETF.

\bibitem{schmitt2012} Schmitt, F., Gassen, J., \& Gerhards-Padilla, E. (2012). HTTP Antivirus Proxy --- An Approach to Defend Against Slow HTTP DoS Attacks. In \textit{Proceedings of IEEE ICCNC 2012}.

\bibitem{alhija2025} Abu Al-Haija, Q., Masoud, Z., Yasin, A., Alesawi, K., \& Alkarnawi, Y. (2025). End-to-End Threat Hunting with a Novel Multiclass Dataset for Intelligent Intrusion Detection. \textit{arXiv preprint arXiv:2508.05609}. Jordan University of Science and Technology \& Princess Sumaya University of Technology.

\end{thebibliography}

\end{document}
